Tato sekce se podrobně zabývá implementací jednotlivých herních mechanik na úrovni kódu. Budou popsány klíčové třídy, jejich metody a algoritmy, které tvoří základ herního systému.

\subsection{PlayerControllerData}

Třída \texttt{PlayerControllerData} představuje klíčový prvek celého systému herních mechanik, neboť definuje veškeré parametry pohybu a chování hráče na jednom centrálním místě. V této sekci bude podrobně rozebrána struktura této třídy, její účel a způsob použití v rámci celého systému.
\newpage
\subsubsection{Účel a využití ScriptableObject}

V Unity je možné vytvořit více druhů skriptů. Nejčastěji používaným je \texttt{MonoBehaviour}, který je určen především pro skripty přidělované na herní objekty v scéně. Pro ukládání dat a konfigurační parametry však není \texttt{MonoBehaviour} ideální volbou, protože vyžaduje přítomnost herního objektu a není možné jej snadno sdílet mezi více scénami.

Z tohoto důvodu existuje typ \texttt{ScriptableObject}, který je navržen speciálně pro ukládání dat a konfiguračních hodnot. \texttt{ScriptableObject} je samostatný datový objekt, který není přidělen k žádnému hernímu objektu a může být vytvořen jako samostatný asset v projektu. Hlavní výhodou tohoto přístupu je možnost mít více různých sad parametrů a snadno je mezi nimi přepínat. Každý \texttt{ScriptableObject} asset může obsahovat odlišné hodnoty, což umožňuje například vytvořit různé profily obtížnosti nebo testovací scénáře.

Pro vytvoření \texttt{ScriptableObject} ve editoru se používá atribut \texttt{[CreateAssetMenu]}, který přidá položku do menu pro vytvoření nového assetu. Tím získá developer možnost tweakovat hodnoty bez nutnosti otevírat samotný skript a měnit v něm konstanty. Při změně hodnot v \texttt{ScriptableObject} assetu se tyto změny okamžitě projeví ve hře, což výrazně zrychluje proces ladění herního pocitu.

\begin{lstlisting}[language=csharp,caption={Listing 1: Definice třídy PlayerControllerData}]
[CreateAssetMenu(menuName = "Player Movement Data")]
public class PlayerControllerData : ScriptableObject
{
    [Header("Run")]
    public float runMaxSpeed = 8f;
    public float runAcceleration = 20f;
    [HideInInspector] public float runAccelAmount;
    
    [Header("Jump")]
    public float jumpForce = 15f;
    public float coyoteTime = 0.1f;
    public float jumpBufferTime = 0.1f;
    
    [Header("Gravity")]
    public float gravity = 9.8f;
    public float fallMultiplier = 2.5f;
    
    [Header("Dash")]
    public float dashSpeed = 20f;
    public float dashTime = 0.15f;
    public float dashCooldown = 0.3f;
    
    [Header("Wall Slide / Jump")]
    public float wallSlideSpeed = 2f;
    public float wallJumpDuration = 0.4f;
    public float wallJumpPowerY = 16f;
}
\end{lstlisting}

\subsection{Struktura dat a organizace}

Třída \texttt{PlayerControllerData} dědí ze \texttt{ScriptableObject} a je organizována do několika sekcí, z nichž každá sdružuje parametry související s konkrétní herní mechanikou. Pro přehledné oddělení jednotlivých sekcí v editoru se používá atribut \texttt{[Header]}, který vytváří vizuální oddělovače mezi skupinami proměnných.

První sekce obsahuje parametry pro běh (run), konkrétně maximální rychlost, zrychlení a poměr zrychlení ve vzduchu. Druhá sekce se zaměřuje na skok a obsahuje sílu skoku, časové parametry pro coyote time a jump buffer. Třetí sekce definuje parametry gravitace, včetně základní gravitace a multiplikátoru pro pád. Čtvrtá sekce obsahuje nastavení pro přerušení skoku (jump cut) a pátá sekce parametry pro dash. Poslední sekce se věnuje wall jump, včetně sil pro jednotlivé typy skoků.

Tato organizace umožňuje rychle nalézt a upravit specifický parametr bez nutnosti procházet neuspořádaný seznam všech proměnných. Zároveň usnadňuje pochopení systému jako celku, protože jednotlivé sekce jasně korespondují s implementovanými herními mechanikami.
\subsection{Důležité parametry a jejich vliv na herní pocit}

Každý parametr v \texttt{PlayerControllerData} má specifický vliv na chování hráče a celkový herní pocit. Pochopení jejich vzájemného působení je klíčové pro správné nastavení mechanik.

Mezi nejdůležitější parametry běhu patří \texttt{runMaxSpeed}, který definuje maximální rychlost horizontálního pohybu. Parametr \texttt{runAcceleration} určuje, jak rychle hráč dosáhne této maximální rychlosti. Vysoká hodnota zrychlení vede k okamžitému startu, zatímco nízká hodnota vytváří efekt rozjezdu. Důležitým parametrem je také \texttt{accelInAir}, který určuje procentuální poměr pozemního zrychlení použitého ve vzduchu. Nízká hodnota tohoto parametru ztěžuje změnu směru během skoku, což může být žádoucí pro obtížnější sekce.

Pro skok jsou klíčové parametry \texttt{jumpForce} určující sílu počátečního výskoku, \texttt{coyoteTime} definující dobu grace period po opuštění platformy, a \texttt{jumpBufferTime} pro registraci skoku před dopadem. Parametr \texttt{jumpHangTimeThreshold} určuje prahovou rychlost, při které se aktivuje efekt vznášení na vrcholu skoku.

Gravitace je řízena parametrem \texttt{gravity} pro základní gravitační sílu a \texttt{fallMultiplier} pro zrychlení pádu. Vyšší hodnota fall multiplieru vede k rychlejšímu pádu a kratšímu času ve vzduchu, což zlepšuje kontrolu nad přesným načasováním dopadu.

Dash je konfigurován pomocí \texttt{dashSpeed} pro rychlost výpadu, \texttt{dashTime} pro dobu trvání a \texttt{dashCooldown} pro prodlevu mezi dashi.
\newpage
\subsubsection{Výpočet akcelerace a její automatizace}

Pro zajištění konzistentního chování pohybu bez ohledu na maximální rychlost se v metodě \texttt{OnValidate()} automaticky počítá hodnota \texttt{runAccelAmount} z parametrů \newline \texttt{runAcceleration} a \texttt{runMaxSpeed}. Tento výpočet je důležitý pro zajištění toho, aby zrychlení bylo proporcionální k maximální rychlosti.

Vztah mezi akcelerací a maximální rychlostí je následující: při nízké maximální rychlosti a vysoké akceleraci by mohl hráč dosáhnout maximální rychlosti příliš rychle, což by vedlo k neintuitivnímu chování. Automatickým výpočtem se zajistí, že pohyb zůstává konzistentní a předvídatelný pro jakoukoli kombinaci hodnot.

Metoda \texttt{OnValidate()} se volá v editoru při změně hodnot v inspektoru, což znamená, že výpočet probíhá automaticky při každé úpravě parametrů. Nemusíme tak manuálně počítat správnou hodnotu akcelerace a může se soustředit pouze na nastavení požadovaného chování.
\subsection{Použití v PlayerController a propojení systémů}

V třídě \texttt{PlayerController} je instance \texttt{PlayerControllerData} přidělena jako serializované pole pomocí atributu \texttt{[SerializeField]}. Toto pole se v editoru zobrazuje jako pole pro přetahování (drag and drop), kam lze umístit libovolný \texttt{ScriptableObject} asset vytvořený z třídy \texttt{PlayerControllerData}.

Při spuštění hry se v metodě \texttt{Awake()} načtou všechny potřebné reference a inicializují se výchozí hodnoty. Parametry z \texttt{PlayerControllerData} se pak používají v celém systému pro řízení chování hráče. Toto propojení umožňuje snadné ladění hodnot bez nutnosti otevírat kód a měnit konstanty přímo ve skriptech.

Výhodou tohoto přístupu je možnost vytvořit více variant \texttt{PlayerControllerData} assetů s různými parametry. Například lze vytvořit jeden asset pro testovací úroveň s nižší obtížností a jiný pro finální úroveň s vyššími nároky na ovládání. Přepínání mezi těmito variantami je pak otázkou jediného přetahovacího pole v inspektoru.

\subsection{PlayerController}

Hlavním řídicím skriptem celého systému pohybu hráče je třída \texttt{PlayerController}. Tento skript funguje jako centrální uzel, který koordinuje všechny ostatní systémy a zajišťuje jejich vzájemnou spolupráci. V této sekci budou popsány jednotlivé aspekty implementace tohoto skriptu, včetně zpracování vstupu, aplikace pohybu a propojení s dalšími systémy.
\newpage
\subsubsection{Hlavní skript a jeho role v systému}

Třída \texttt{PlayerController} představuje hlavní skript propojující všechny systémy pohybu hráče. Dědí z \texttt{MonoBehaviour} a přidává se na herní objekt hráče jako komponenta. Tento skript funguje jako centrální uzel, který koordinuje vstupy, fyziku, gravitaci a propojení s dalšími systémy jako dash, střelba a wall jump.

V metodě \texttt{Awake()} probíhá inicializace skriptu. Získávají se reference na komponentu \texttt{Rigidbody2D}, která je esenciální pro fyzikální pohyb, a na ostatní skripty jako \texttt{PlayerDash}, \texttt{PlayerStripeShot} a \texttt{PlayerWall}. Důležitým krokem je nastavení \newline \texttt{gravityScale} na hodnotu 0 na komponentě \texttt{Rigidbody2D}. To je klíčové pro fungování vlastního gravitačního systému, protože Unity vestavěná gravitace je kompletně deaktivována a veškerá gravitace je řešena vlastním kódem v metodě \texttt{ApplyCustomGravity()}.

\begin{lstlisting}[language=csharp]
private void Awake()
{
    rb = GetComponent<Rigidbody2D>();
    rb.gravityScale = 0;
    isFacingRight = true;
    pd = GetComponent<PlayerDash>();
    pss = GetComponent<PlayerStripeShot>();
    pw = GetComponent<PlayerWall>();
}
\end{lstlisting}

\newpage
\subsubsection{Input systém a zpracování uživatelského vstupu}

Vstupy od uživatele se zpracovávají pomocí Unity Input Systemu, který poskytuje modernější a flexibilnější přístup ke zpracování vstupů oproti staršímu Input Manageru. Pomocí atributu \texttt{[SerializeField] private InputActionReference} se přiřazují jednotlivé akce v editoru. Každá akce repreguje konkrétní typ vstupu, jako je pohyb, skok nebo střelba.

K registraci inputu dochází v metodě \texttt{OnEnable()}, kde se připojují callbacky na jednotlivé akce. Tyto callbacky jsou metody, které se automaticky volají při změně stavu vstupu. Příkladem je \texttt{OnMoveInput()} pro pohyb, \texttt{OnJumpInput()} pro skok a podobně. Při uvolnění skriptu nebo deaktivaci objektu se v \texttt{OnDisable()} tyto callbacky odpojují, což zabraňuje chybám při volání metod na neexistující objekty.

Tento přístup k input systému je výhodný, protože umožňuje snadnou rekonfiguraci vstupů bez nutnosti měnit kód. V editoru lze přiřadit libovolnou klávesu nebo tlačítko k akci a systém automaticky zpracuje změnu.

\begin{lstlisting}[language=csharp,caption={Listing 3: Deklarace vstupních akcí}]
[Header("input")]
public Vector2 moveInput { get; private set; }
[SerializeField] private InputActionReference moveAction;
[SerializeField] private InputActionReference jumpAction;
[SerializeField] private InputActionReference dashAction;
[SerializeField] private InputActionReference shootAction;

private void OnEnable()
{
    moveAction.action.Enable();
    jumpAction.action.Enable();
    shootAction.action.Enable();
    dashAction.action.Enable();

    jumpAction.action.performed += OnJumpInput;
    jumpAction.action.canceled += OnJumpRelease;
    dashAction.action.performed += OnDashInput;
    shootAction.action.performed += OnShootInput;
}
\end{lstlisting}

\newpage
\subsection{Zpracování vstupu a filtrování hodnot}

V metodě \texttt{Update()} se čte vstup z akce pohybu a následně filtruje pomocí prahu 0.1. Tento práh slouží k ignorování nechtěného inputu z joysticku, který může generovat velmi malé hodnoty i při neutrální pozici. Bez tohoto filtru by se hráč mohl nechtěně pohybovat i když joystick není aktivně ovládán. Tato logika platí také pro input z klávesnice.

Hodnota vstupu se ukládá jako \texttt{Vector2}, který obsahuje složky x a y. Podle hodnoty na ose x se následně volá metoda \texttt{CheckDirectionToFace()}, která kontroluje otočení postavy. Pokud je vstup kladný a postava se aktuálně dívala doleva, postava se otočí doprava. Podobná logika funguje také v opačném scénáři. Toto zajišťuje, že postava vždy směřuje ve směru pohybu.

Metoda \texttt{Update()} je vhodná pro zpracování vstupu, protože běží jednou za snímek a není vázána na fyzikální timestep. To umožňuje responsivní zpracování inputu bez ohledu na fyzikální nastavení hry.
\newpage
\subsection{Horizontální pohyb a akcelerace}

Metoda \texttt{Run()} aplikuje horizontální pohyb pomocí \texttt{Lerp} interpolace, která zajišťuje plynulé zrychlení a zpomalování namísto okamžité změny rychlosti. Vzorec \texttt{Mathf.Lerp(rb.linearVelocity.x, targetSpeed, Time.fixedDeltaTime * accelRate)} postupně přesouvá aktuální rychlost k cílové rychlosti.

Zrychlení se liší podle toho, zda je hráč na zemi nebo ve vzduchu. Parametr \texttt{accelInAir} určuje, kolik procent pozemního zrychlení se použije ve vzduchu. Typicky se používá nižší hodnota (například 50-70 procent), což zajišťuje, že hráč nemůže ve vzduchu měnit směr tak rychle jako na zemi. Toto omezení přidává hloubku do herní mechaniky a odměňuje hráče za správné načasování pohybu.

Během provádění wall jump je volání \texttt{Run()} dočasně deaktivováno, aby hráč nemohl během letu po stěně reagovat na směrové tlačítko. Toto zajišťuje, že wall jump má konzistentní trajektorii bez ohledu na input hráče během letu.

\begin{lstlisting}[language=csharp,caption={Listing 4: Metoda Run pro horizontální pohyb}]
private void Run()
{
    if (pw != null && pw.IsWallJumping)
        return;
    if (rb == null || data == null) return;

    float targetSpeed = moveInput.x * data.runMaxSpeed;

    float accelRate = lastOnGroundTime > 0 ?
        data.runAccelAmount :
        data.runAccelAmount * data.accelInAir;

    if (isJumping && Mathf.Abs(rb.linearVelocity.y) < data.jumpHangTimeThreshold)
    {
        accelRate *= data.jumpHangAccelerationMult;
        targetSpeed *= data.jumpHangMaxSpeedMult;
    }

    float newVelX = Mathf.Lerp(rb.linearVelocity.x, targetSpeed, 
        Time.fixedDeltaTime * accelRate);
    rb.linearVelocity = new Vector2(newVelX, rb.linearVelocity.y);
}
\end{lstlisting}

\newpage
\subsection{Vlastní gravitace a její dynamické řízení}

Jelikož je \texttt{gravityScale} na \texttt{Rigidbody} nastaven na 0, gravitace se kompletně řeší vlastním kódem v metodě \texttt{ApplyCustomGravity()}. Toto řešení poskytuje plnou kontrolu nad chováním hráče ve vzduchu a umožňuje dynamické úpravy gravitace podle aktuálního stavu.

Metoda rozlišuje čtyři různé stavy, z nichž každý má specifickou gravitační sílu. Prvním stavem je normální pád dolů, při kterém se použije \texttt{fallMultiplier} pro rychlejší pád. To zajistí, že hráč padá rychleji, než stoupá, což je typické pro kvalitní platformovky. Druhým stavem je přerušení skoku (jump cut), kdy hráč pustí tlačítko skoku během vzestupu. V tomto případě se použije ještě vyšší multiplikátor gravitace pro okamžité ukončení vzestupu. Třetím stavem je dosažení vrcholu skoku, kdy vertikální rychlost klesne pod \texttt{jumpHangTimeThreshold}. V tomto momentě se gravitace dočasně sníží, což vytváří efekt krátkého vznášení. Posledním stavem je normální stav bez jakékoli modifikace, kdy se použije základní gravitace.

\begin{lstlisting}[language=csharp,caption={Listing 5: Metoda ApplyCustomGravity}]
private void ApplyCustomGravity()
{
    if (rb == null || data == null) return;

    Vector2 vel = rb.linearVelocity;
    float g = Physics2D.gravity.y * data.gravity * Time.fixedDeltaTime;

    if (vel.y < 0)
        vel.y += g * data.fallMultiplier;
    else if (jumpCut)
        vel.y += g * data.jumpCutMultiplier;
    else if (isJumping && Mathf.Abs(vel.y) < data.jumpHangTimeThreshold)
        vel.y += g * data.jumpHangGravityMult;
    else
        vel.y += g;

    rb.linearVelocity = vel;
}
\end{lstlisting}

\subsection{Výpočet nové rychlosti a použití Lerp}

Při aplikaci pohybu se používá interpolace pomocí \texttt{Lerp} pro zajištění plynulého zrychlení a zpomalování. Tento přístup je výhodnější než přímé nastavení rychlosti, protože vytváří přirozenější herní pocit. Hráč nedosáhne maximální rychlosti okamžitě, ale postupně se k ní přibližuje.

Vzorec \texttt{Mathf.Lerp(rb.linearVelocity.x, targetSpeed, Time.fixedDeltaTime \newline * accelRate)} pracuje následovně: první parametr je aktuální rychlost na ose x, druhý parametr je cílová rychlost (buď maximální rychlost ve směru pohybu nebo 0 při zastavení) a třetí parametr určuje rychlost interpolace. Parametr \texttt{Time.fixedDeltaTime} zajišťuje, že interpolace je nezávislá na snímkové frekvenci, zatímco \texttt{accelRate} určuje, jak rychle se má hráč accelerovat.

Hodnota \texttt{targetSpeed} se vypočítá vynásobením vstupu (1 nebo -1) s maximální rychlostí. To zajišťuje, že hráč běží maximální rychlostí, pokud drží směrové tlačítko, a zastaví se, pokud žádné tlačítko nedrží.
\subsection{Zemní kontrola a detekce dotyku se zemí}

Kontrola, zda se hráč dotýká země, je fundamentální pro správné fungování skoku a dalších mechanik. V \texttt{PlayerController} se tato kontrola provádí metodou \texttt{IsGrounded()}, která využívá \texttt{Physics2D.OverlapBox}.

Metoda \texttt{Physics2D.OverlapBox} vytváří virtuální box na specifikované pozici a kontroluje, zda se v tomto prostoru nachází nějaký collider. Parametry zahrnují \newline \texttt{groundCheckPoint} udávající střed boxu a \texttt{groundCheckSize} udávající rozměry boxu. Box je typicky malý a umístěný na spodní část nohou hráče.

Výsledek detekce se ukládá do proměnné \texttt{grounded}, která se následně používá pro rozhodování, zda je možné provést skok, jakou akceleraci použít a další logické operace. Detekce probíhá v každém snímku, což zajišťuje aktuální stav bez zpoždění.
\subsubsection{Propojení s Animatorem a předávání parametrů}

Animace jsou důležitou součástí herního pocitu a přispívají k celkovému dojmu ze hry. V metodě \texttt{Update()} se předávají hodnoty do Animatoru, který následně přepíná mezi jednotlivými stavy animací podle aktuálního stavu hráče.

Prvním předávaným parametrem je rychlost na ose x (\texttt{speed}), která se používá pro animaci běhu. Hodnota se typicky pohybuje od 0 do 1 a Animator ji používá k přechodu mezi stavy idle a run, případně k nastavení rychlosti animace.

Druhým parametrem je \texttt{isGrounded}, což je boolean hodnota indikující, zda se hráč dotýká země. Tento parametr se používá pro přepínání mezi animací stání a animací vzduchu (skok, pád). Třetím parametrem je vertikální rychlost (\texttt{verticalVelocity}), která slouží k detekci pádu. Pokud je vertikální rychlost záporná a hráč není na zemi, spustí se animace pádu.

Správné propojení mezi logikou a animacemi je klíčové pro plynulý herní zážitek a eliminaci vizuálních nesrovnalostí mezi chováním postavy a její animací.
\subsubsection{Otočení postavy a vizuální změna směru}

Metoda \texttt{CheckDirectionToFace()} kontroluje, zda se postava má otočit, a provádí vizuální změnu směru. Na rozdíl od fyzikální změny směru, která by vyžadovala změnu vektoru rychlosti, tato metoda pouze mění vizuální prezentaci postavy.

Otočení se provádí překlopením \texttt{localScale.x} na zápornou hodnotu. Pokud je postava otočená doprava a vstup je doleva (záporná hodnota), \texttt{localScale.x} se nastaví na -1. Podobně při otočení doleva a vstupu doprava se \texttt{localScale.x} nastaví na 1.

Tento přístup je výhodný, protože nevyžaduje změnu fyzikálního směru, který je řízen rychlostí na \texttt{Rigidbody2D}. Postava vizuálně směřuje správným směrem, zatímco fyzika může být řešena nezávisle. To je užitečné například při střelbě, kde projektil vždy směřuje k cíli bez ohledu na vizuální orientaci postavy.
\subsubsection{Propojení s dash systémem a deaktivace pohybu}

V metodě \texttt{FixedUpdate()} třídy \texttt{PlayerController} se kontroluje, zda je aktivní dash. Pokud ano, přeskočí se volání \texttt{Run()} a \texttt{ApplyCustomGravity()}. Toto zajišťuje, že během dashe hráč nereaguje na gravitaci ani na vstup pohybu.

Důvod pro deaktivaci gravitace během dashe je ten, že dash by měl být konzistentní bez ohledu na vertikální polohu. Pokud by gravitace byla stále aktivní, hráč by během dashe klesal, což by narušilo dojem rychlého výpadu. Deaktivace vstupu pohybu zajišťuje, že dash má předem definovanou trajektorii a hráč ji nemůže během dashe měnit.

Propojení mezi \texttt{PlayerController} a \texttt{PlayerDash} je implementováno pomocí volání metody \texttt{IsDashing()} v \texttt{PlayerDash}, která vrací aktuální stav dashe. Toto rozhodnutí se pak používá v \texttt{FixedUpdate()} pro přeskočení standardních herních mechanik.

\subsection{Skok}

Systém skoku představuje jednu z nejdůležitějších součástí každé 2D plošinovky a jeho správná implementace zásadně ovlivňuje celkový herní pocit. V této sekci budou podrobně rozebrány všechny implementované mechaniky související se skokem, včetně pokročilých technik jako je coyote time, jump buffer a proměnná výška skoku.

\newpage
\subsubsection{Coyote time a jeho implementace}

Jak už bylo zmíněno v koncepci, coyote time dovolí hráči skočit i pár snímků poté, co opustil platformu. V \texttt{PlayerController} se to řeší pomocí proměnné \texttt{lastOnGroundTime}, která se při dopadu na zem nastaví na hodnotu \texttt{coyoteTime} z \texttt{PlayerControllerData}. Následně se v každém snímku v metodě \texttt{Update()} tato proměnná snižuje o \texttt{Time.deltaTime}.

V momentě, kdy hráč stiskne skok, se kontroluje, zda je \texttt{lastOnGroundTime} větší než 0. Pokud ano, skok se provede i když hráč technicky už není na zemi. Toto řešení poskytuje hráči malé okno neperfektnosti, kdy se skok stále zaregistruje, i když už by hráč měl správně spadnout.

Hodnota \texttt{coyoteTime} se typicky nastavuje na 0.1 až 0.2 sekundy, což odpovídá přibližně 6-12 snímkům při 60 snímcích za sekundu. Tato doba je dostatečně krátká, aby neovlivnila běžnou hru, ale dostatečně dlouhá, aby zachytila nepatrně opožděný input hráče.

\begin{lstlisting}[language=csharp,caption={Listing 6: Implementace coyote time}]
// In Update() the timer decreases
private void Update()
{
    lastOnGroundTime -= Time.deltaTime;
    lastPressedJumpTime -= Time.deltaTime;
    IsGrounded();
    
    // Set coyote time when grounded
    if (grounded) 
        lastOnGroundTime = data.coyoteTime;
    
    // Check jump with coyote time
    if (lastPressedJumpTime > 0 && lastOnGroundTime > 0 && !pw.IsWallJumping)
        Jump();
}
\end{lstlisting}

\subsection{Jump buffer a registrace skoku před dopadem}

Jump buffer funguje obdobně jako coyote time, ale na opačnou stranu časové osy. Zatímco coyote time řeší situaci po opuštění platformy, jump buffer řeší situaci před dopadem na zem. Když hra zaznamená input od hráče v momentě, kdy hráč teprve dopadá z výšky na zem, ovšem \texttt{groundCheck} ještě nezaznamenal \texttt{ground layer}, avšak je pár snímků nad ní, tento input se zaregistruje jako skok.

Implementace využívá proměnnou \texttt{lastPressedJumpTime}, která se nastaví při stisknutí skoku na hodnotu \texttt{jumpBufferTime} z \texttt{PlayerControllerData}. V každém snímku se tato hodnota snižuje o \texttt{Time.deltaTime}.

V \texttt{Update()} se pak kontroluje, zda je \texttt{lastPressedJumpTime > 0} a současně \texttt{lastOnGroundTime > 0}. Pokud obě podmínky platí, provede se skok. Toto řešení eliminuje frustraci z toho, že hráč stiskne skok těsně před dopadem a skok se neprovede, protože hráč ještě není na zemi.
\subsubsection{Provádění skoku a nastavení fyzikálních parametrů}

Metoda \texttt{Jump()} se volá, když jsou splněny všechny podmínky pro skok (dostatečný coyote time a jump buffer). Nejprve se nastaví \texttt{isJumping} na true, což indikuje, že hráč je aktuálně ve vzduchu a nelze skočit znovu. Současně se nastaví \texttt{jumpCut} na false, což resetuje stav přerušení skoku.

Následně se vynuluje vertikální rychlost na \texttt{Rigidbody2D} a poté se nastaví na hodnotu \texttt{jumpForce} z \texttt{PlayerControllerData}. Toto zajistí konzistentní počáteční rychlost skoku bez ohledu na předchozí pohyb. Důležité je také vynulování časových proměnných \texttt{lastOnGroundTime} a \texttt{lastPressedJumpTime}, aby se skok neopakoval v dalších snímcích.

Hodnota \texttt{jumpForce} se typicky pohybuje v rozmezí 10-20 v závislosti na nastavení gravitace a požadované výšce skoku. Správné nastavení této hodnoty je klíčové pro celkový herní pocit.

\begin{lstlisting}[language=csharp,caption={Listing 7: Metoda Jump pro provedení skoku}]
private void Jump()
{
    if (rb == null || data == null) return;

    isJumping = true;
    jumpCut = false;
    lastPressedJumpTime = 0;
    lastOnGroundTime = 0;

    rb.linearVelocity = new Vector2(rb.linearVelocity.x, 0);
    rb.linearVelocity = new Vector2(rb.linearVelocity.x, data.jumpForce);
}
\end{lstlisting}

\newpage
\subsection{Jump cut a předčasné ukončení skoku}

Při puštění tlačítka skoku během vzestupu se aktivuje jump cut, který předčasně ukončí skok a umožní hráči rychleji spadnout. Toto je důležitá mechanika, která umožňuje hráči kontrolovat výšku skoku podle délky stisknutí tlačítka.

Implementace využívá metodu \texttt{OnJumpRelease()}, která se registruje jako callback při uvolnění vstupu skoku. V této metodě se kontroluje, zda je hráčova vertikální rychlost kladná (tedy stoupá). Pokud ano, vynásobí se rychlost hodnotou \texttt{jumpCutVelocityMultiplier} z \texttt{PlayerControllerData}. Typicky se používá hodnota 0.5, což znamená, že při puštění tlačítka se rychlost sníží na polovinu.

V metodě \texttt{ApplyCustomGravity()} se pak při aktivním \texttt{jumpCut} aplikuje ještě vyšší gravitace než při normálním pádu. Toto zajistí okamžité ukončení vzestupu a rychlý přechod do pádu, což přispívá k responsivitě ovládání.

\begin{lstlisting}[language=csharp,caption={Listing 8: Metoda OnJumpRelease pro jump cut}]
private void OnJumpRelease(InputAction.CallbackContext ctx)
{
    if (rb != null && rb.linearVelocity.y > 0)
    {
        jumpCut = true;

        rb.linearVelocity = new Vector2(
            rb.linearVelocity.x,
            rb.linearVelocity.y * data.jumpCutVelocityMultiplier
        );
    }
}
\end{lstlisting}

\subsubsection{Jump hang a efekt vznášení na vrcholu skoku}

Po dosažení vrcholu skoku, kdy je vertikální rychlost nižší než \texttt{jumpHangTimeThreshold} z \texttt{PlayerControllerData}, se dočasně sníží gravitace. Tím vzniká krátký moment vznášení, který výrazně zlepšuje herní pocit a poskytuje hráči více času na načasování dalších akcí.

Implementace probíhá v metodě \texttt{ApplyCustomGravity()}, kde se kontroluje, zda je aktivní skok a zda je vertikální rychlost nižší než práh. Pokud ano, použije se snížená gravitace. Zároveň se v metodě \texttt{Run()} detekuje tento stav a mírně se zvýší akcelerace a maximální rychlost.

Toto zvýšení mobility během jump hang poskytuje hráči lepší kontrolu nad pohybem ve vzduchu v blízkosti vrcholu skoku. Hráč může lépe reagovat na herní situace a přesněji se dostat na cílovou platformu.
\newpage
\subsection{Vstup skoku a registrace callbacků}

Skok se registruje jako input callback v metodě \texttt{OnJumpInput()}, která se připojí k Unity Input System akci skoku. V této metodě se nastavuje \texttt{lastPressedJumpTime} na hodnotu \texttt{jumpBufferTime}, což aktivuje jump buffer systém.

Proměnná \texttt{jumpPressedThisFrame} se nastaví na true, což indikuje, že v aktuálním snímku byl stisknut skok. Tato proměnná se používá pro okamžitou detekci stisknutí v rámci snímku. Na konci \texttt{Update()} se \texttt{jumpPressedThisFrame} vynuluje, aby se zabránilo opakovanému skoku v dalších snímcích.

Při uvolnění tlačítka skoku se volá \texttt{OnJumpRelease()}, která aktivuje \texttt{jumpCut}. Toto rozdělení na stisk a uvolnění umožňuje jemnou kontrolu nad výškou skoku a přispívá k celkové responsivitě ovládání.

\subsection{Dash}

Mechanika rychlého výpadu neboli dash představuje další důležitý prvek pohybového systému. Dash umožňuje hráči na krátký časový okamžik rapidní zrychlení a přesun do aktuálního směru pohybu. Tato mechanika při správném nastavení výrazně zvyšuje plynulost hry a celkově zpříjemňuje pocit z pohybu.

\subsubsection{Základní princip a dědění z PlayerAbility}

Mechanika rychlého výpadu je implementována ve třídě \texttt{PlayerDash}, která dědí ze třídy \texttt{PlayerAbility}. \texttt{PlayerAbility} je abstraktní třída, která poskytuje základní systém pro odemykání a zamykání jednotlivých schopností. Všechny ability skripty v projektu dědí z této abstraktní třídy.

Třída \texttt{PlayerAbility} obsahuje pouze proměnnou \texttt{unlocked} typu boolean a metody \texttt{Unlock()} a \texttt{Lock()}. Při volání \texttt{Unlock()} se \texttt{unlocked} nastaví na true a schopnost je k dispozici. Při volání \texttt{Lock()} se schopnost deaktivuje a hráč ji nemůže používat.

Tato architektura umožňuje snadné přidávání nových schopností bez nutnosti opakovat kód pro odemykání a zamykání. Každá konkrétní schopnost pouze implementuje svou vlastní logiku, zatímco systém odemykání je společný.
\newpage
\subsubsection{Metoda TryDash a kontrola podmínek}

Metoda \texttt{TryDash()} je veřejná metoda volaná z \texttt{PlayerController} při inputu dashe. Nejprve kontroluje, zda je schopnost odemčená voláním \texttt{IsUnlocked()}. Pokud není, metoda se ukončí a dash se neprovede.

Následně se kontroluje cooldown a stav dashe. Pokud běží cooldown \newline (\texttt{dashCooldownTimer > 0}) nebo pokud je dash aktuálně aktivní, metoda se ukončí. Teprve pokud všechny tyto podmínky jsou splněny, pokračuje se k samotnému provedení dashe.

Po úspěšné kontrole podmínek se určí směr dashe. Pokud hráč dává vstup pohybu (\texttt{moveInput} není nulový), použije se tento směr. V opačném případě se použije směr podle otočení postavy jako záloha, aby dash vždy fungoval i když hráč stojí na místě.

\begin{lstlisting}[language=csharp,caption={Listing 9: Metoda TryDash}]
public void TryDash()
{
    if (!unlocked) return;

    if (isDashing || dashCooldownTimer > 0 || rb == null || data == null) return;

    Vector2 moveInput = pc != null ? pc.moveInput : Vector2.zero;
    dashDir = new Vector2(moveInput.x, moveInput.y).normalized;
    if (dashDir == Vector2.zero)
        dashDir = pc.isFacingRight ? Vector2.right : Vector2.left;

    isDashing = true;
    dashTimer = data.dashTime;
    dashCooldownTimer = data.dashCooldown;
    rb.linearVelocity = dashDir * data.dashSpeed;
}
\end{lstlisting}

\subsection{Směr dashe a normalizace vektoru}

Směr dashe se normalizuje, aby byl vždy jednotkový vektor. To zajišťuje, že dash má konzistentní rychlost bez ohledu na úhel vstupu. Pokud by normalizace nebyla provedena, diagonální dash by byl rychlejší než horizontální nebo vertikální.

Normalizace se provádí metodou \texttt{Normalize()} na \texttt{Vector2}. Pokud je vektor příliš krátký (méně než 0.01), považuje se za nulový a nepokračuje se v dashi.

V případě, že hráč nedává žádný vstup pohybu, použije se směr podle otočení postavy. To se zjistí pomocí proměnné \texttt{pc.isFacingRight}, která vrací true pokud se postava dívá doprava. Výsledný směr je pak \texttt{Vector2.right} nebo \texttt{Vector2.left}. Toto zajišťuje, že dash vždy funguje, i když hráč stojí na místě a nedává žádný vstup.
\subsection{DashMovement a aplikace rychlosti během dashe}

Po dobu trvání dashe se ve \texttt{FixedUpdate()} přepisuje rychlost \texttt{Rigidbody} na konstantní hodnotu \texttt{dashDir * dashSpeed}. Toto zajišťuje, že hráč se pohybuje konstantní rychlostí bez ohledu na vstup nebo gravitaci. Vzorec je \texttt{rb.linearVelocity = \newline new Vector2(dashDir.x * dashSpeed, dashDir.y * dashSpeed)}.

Po uplynutí \texttt{dashTime} se dash ukončí voláním \texttt{EndDash()}. Zároveň se nastaví cooldown timer na hodnotu \texttt{dashCooldown}. Po celou dobu trvání dashe se \texttt{Rigidbody} chová jako by měl nulovou hmotnost, protože rychlost je přepisována každý snímek. To znamená, že kolize během dashe fungují správně, ale hráč je řízen čistě logikou dashe.

\begin{lstlisting}[language=CSharp,caption={Listing 10: Metoda DashMovement}]
private void FixedUpdate()
{
    if (isDashing)
    {
        DashMovement();
    }
}

private void DashMovement()
{
    if (rb == null || data == null) return;

    dashTimer -= Time.fixedDeltaTime;
    rb.linearVelocity = dashDir * data.dashSpeed;

    if (dashTimer <= 0)
    {
        isDashing = false;
    }
}

public bool IsDashing()
{
    return isDashing;
}
\end{lstlisting}
\newpage
\subsection{Cooldown systém a omezení frekvence dashe}

Po dokončení dashe se nastaví \texttt{dashCooldownTimer} na hodnotu \texttt{dashCooldown} z \newline \texttt{PlayerControllerData}. V \texttt{Update()} se tento timer snižuje o \texttt{Time.deltaTime}. Dash lze provést pouze když je timer menší nebo roven nule.

Tento cooldown systém zabraňuje spamování dashe a vyvažuje herní mechaniku. Bez cooldownu by hráč mohl dash používat téměř neustále, což by změnilo charakter hry. Správně nastavený cooldown nutí hráče strategicky využívat dash a rozhodovat se, kdy je nejvhodnější ho použít.

Hodnota \texttt{dashCooldown} se typicky nastavuje na 0.5-2 sekundy v závislosti na \newline požadované obtížnosti a délce úrovní. Kratší cooldown umožňuje agresivnější hru, zatímco delší cooldown vyžaduje opatrnější přístup.
\subsection{Propojení s PlayerController a registrace inputu}

Ve \texttt{FixedUpdate()} třídy \texttt{PlayerController} se při aktivním dashi přeskočí volání \texttt{Run()} a \texttt{ApplyCustomGravity()}, jak bylo popsáno v předchozí sekci. Toto propojení zajišťuje, že všechny herní systémy spolupracují a nedochází ke konfliktům.

Dash se registruje jako input callback v metodě \texttt{OnDashInput()}, kde se volá \newline \texttt{pd.TryDash()}. \texttt{pd} je reference na instanci \texttt{PlayerDash}, která se získá v \texttt{Awake()}. Metoda \texttt{IsDashing()} vrací aktuální stav dashe pro externí skripty, které potřebují vědět, zda hráč právě dashuje.

Toto propojení umožňuje čistou separaci kódu. \texttt{PlayerController} se stará o základní pohyb a gravitaci, zatímco \texttt{PlayerDash} se stará o specifickou logiku dashe. Spolupráce mezi těmito systémy je řízena přes veřejné metody a vlastnosti.

\subsection{Wall jump a wall slide}

Další implementovanou mechanikou je skákání a klouzání po stěnách, tedy wall jump a wall slide. Tyto dvě funkce spolu úzce souvisejí a výrazně rozšiřují možnosti navigace v herním světě. Wall jump umožňuje hráči opětovně vyskočit při kontaktu se stěnou ze strany, čímž se otevírá možnost postupného pohybu mezi dvěma vertikálními plochami a dosažení vyšších částí úrovně.

\subsubsection{Základ skriptu a dědění z PlayerAbility}

Skript \texttt{PlayerWall.cs} dědí z \texttt{PlayerAbility} stejně jako \texttt{PlayerDash}.

Kontrola stěny se provádí pomocí \texttt{Physics2D.OverlapCircle} s bodem \texttt{wallCheck} a poloměrem \texttt{wallCheckRadius} na příslušné vrstvě (layer). Tato metoda detekuje, zda se v okolí bodu \texttt{wallCheck} nachází collider patřící do definované vrstvy. Výsledek určuje, zda je hráč v kontaktu se stěnou.

Podobně jako u \textit{coyote time} se používá \texttt{lastOnWallTime}, který udržuje krátký časový window po opuštění stěny. Toto umožňuje provést wall jump i krátce poté, co hráč technicky opustil stěnu, což zlepšuje herní pocit a snižuje frustrující situace.
\subsubsection{Detekce stěny a určení jejího směru}

V \texttt{Update()} se detekuje, zda se hráč dotýká stěny. Pokud ano a hráč není na zemi, nastaví se \texttt{lastOnWallTime} na \texttt{wallCoyoteTime}. Pokud ne, timer se snižuje o \texttt{Time.deltaTime}. Toto zajišťuje, že hráč má krátký čas na provedení wall jump poté, co opustil stěnu.

Metoda \texttt{DetectWallDirection()} určuje, na které straně stěna je. Existují dva přístupy: buď podle vstupu hráče (pokud hráč drží tlačítko směrem ke stěně, stěna je na této straně), nebo podle otočení postavy jako záloha. První přístup je přesnější, protože umožňuje hráči určit stěnu i když se od ní vzdaluje, druhý přístup slouží jako fallback pro situace bez inputu.
\subsubsection{Wall slide a zpomalený pád podél stěny}

V metodě \texttt{HandleWallSlide()} se kontroluje několik podmínek současně: hráč nesmí být na zemi, musí být na stěně, musí padat dolů (negativní vertikální rychlost), a musí držet pohyb směrem ke stěně. Pokud všechny podmínky platí, vertikální rychlost se ořízne na maximální hodnotu \texttt{wallSlideSpeed}. Tím vzniká efekt zpomaleného klouzání dolů.

Bez této úpravy by hráč podléhal standardní gravitaci a rychleji by padal podél stěny. \texttt{Wall slide} upravuje vertikální rychlost a poskytuje hráči časový prostor pro provedení dalšího skoku. Důležité je, že klouzání se deaktivuje během wall jump, aby nenarušovalo trajektorii odrazu.

Hodnota \texttt{wallSlideSpeed} se typicky nastavuje na nízkou hodnotu, což umožňuje pomalé klouzání bez úplného zastavení. To zajišťuje, že hráč stále padá i na stěně, ale v kontrolovaném tempu.
\newpage
\subsection{Wall jump a odraz od stěny}

Wall jump se provádí pouze když hráč stiskne skok, nachází se na stěně a není na zemi. Skript detekuje směr stěny pomocí \texttt{DetectWallDirection()} a aplikuje sílu opačným směrem. Vertikální složka skoku je pevně nastavena na \texttt{wallJumpPowerY}, zatímco horizontální složka závisí na typu skoku.

Implementace využívá \texttt{Rigidbody2D} pro aplikaci síly metodou \texttt{AddForce()}. Vektory se nastavují podle směru stěny: pokud je stěna vpravo, horizontální složka je záporná, pokud je stěna vlevo, horizontální složka je kladná.

Tento mechanismus umožňuje hráči systematicky se odrážet mezi dvěma stěnami a dosahovat vyšších částí úrovně, kam by se normálním skokem nedostal. Wall jump je základní technikou v mnoha platformovkách a přidává hloubku do navigace herním světem.

\begin{lstlisting}[language=csharp,caption={Listing 11: Metoda HandleWallJump}]
void HandleWallJump()
{
    if (!pc.jumpPressedThisFrame || lastOnWallTime <= 0 || pc.grounded) return;

    isWallJumping = true;
    int jumpDir = -wallDir;

    bool isWallToWall = Physics2D.Raycast(wallCheck.position, 
        Vector2.right * jumpDir, 1f, wallLayer).collider == null;

    float xForce = isWallToWall ? data.wallToWallJumpForceX : data.singleWallJumpForceX;
    rb.linearVelocity = new Vector2(jumpDir * xForce, data.wallJumpPowerY);

    if (pc.isFacingRight != (jumpDir > 0))
        pc.CheckDirectionToFace(jumpDir > 0);

    lastOnWallTime = 0;
}
\end{lstlisting}

\newpage
\subsubsection{Wall-to-wall jump a rychlý přeskok mezi stěnami}

Zajímavou funkcí je rozpoznání \textit{wall-to-wall} skoku. Pomocí \texttt{Raycast} se detekuje, zda je na druhé straně hráče další stěna v dosahu. Pokud ano, použije se vyšší horizontální síla \texttt{wallToWallJumpForceX}, což umožňuje rychlejší přeskok mezi dvěma stěnami. Pokud ne, použije se nižší \texttt{singleWallJumpForceX}.

Rozdíl mezi těmito dvěma silami je významný. Wall-to-wall jump potřebuje více horizontálního impulzu, aby hráč překonal vzdálenost mezi stěnami v krátkém čase. Single wall jump má nižší horizontální složku, jelikož cílem v tomto případě je hráče jen mírně odrazit od stěny pro vizuálně přívětivější efekt ve hře.

Tato funkce přidává do hry další úroveň kontroly a umožňuje zkušeným hráčům rychleji navigovat mezi stěnami. Rozpoznání wall-to-wall situace probíhá vždy před provedením wall jump, takže hráč nemusí manuálně přepínat mezi režimy.
\subsection{Propojení s PlayerController a deaktivace vstupu}

Po provedení wall jump se nastaví \texttt{isWallJumping} na true a odpočítává se timer \newline \texttt{wallJumpTimer}. Dokud \texttt{isWallJumping} trvá, metoda \texttt{Run()} v \texttt{PlayerController} ignoruje horizontální vstup, takže hráč nemůže během wall jump reagovat na směrové tlačítko.

Po uplynutí timeru se \texttt{isWallJumping} vynuluje ve \texttt{FixedUpdate()}. Toto zajišťuje, že wall jump má předem definovanou trajektorii a hráč ji nemůže během letu měnit, což by mohlo narušit konzistenci mechaniky.

Propojení mezi \texttt{PlayerWall} a \texttt{PlayerController} je implementováno přes veřejné metody a vlastnosti. \texttt{PlayerWall} poskytuje metody jako \texttt{CanWallSlide()}, \texttt{TryWallJump()} a vlastnost \texttt{IsWallJumping}. \texttt{PlayerController} tyto metody volá v průběhu herního cyklu.

\subsection{Střelba}

Systém střelby přidává do hry další rozměr interaktivity a strategie. Na rozdíl od tradičních stříleček, kde je střelba primárním prostředkem eliminace nepřátel, v této hře slouží střelba jako záchranný mechanismus. Hráč má k dispozici pouze dva náboje, které se přebíjejí až po vystřelení obou a po krátkém časovém intervalu. Tímto designem podporujeme myšlenku, aby střelba byla pouze poslední možností záchrany a hráč se musel především soustředit na zlepšování svých dovedností v pohybu.

\subsection{Systém střelby a architektura skriptů}

Systém střelby se skládá ze dvou částí -- skriptu na hráči \texttt{PlayerStripeShot.cs} a skriptu na projektilu \texttt{BulletStripeShot.cs}. Toto rozdělení je výhodné pro optimalizaci, protože skript projektilu se aktivuje až při výstřelu a není celou dobu aktivní v paměti.

Hráčský skript \texttt{PlayerStripeShot.cs} dědí z \texttt{PlayerAbility}, což umožňuje odemykat střelbu v průběhu hry jako speciální schopnost. Instance projektilu se provádí až při výstřelu, což šetří výkon. Mezi výstřely skript pouze čeká na input a spravuje stav munice.

Projektily jsou vytvořeny jako prefaby v editoru a obsahují komponenty \texttt{Rigidbody2D} pro fyzikální pohyb a \texttt{Collider2D} pro detekci kolizí. Prefab také obsahuje vizuální komponenty jako sprite renderer pro grafické zobrazení projektilu.
\subsubsection{Směr střelby a automatické zaměření na cíl}

Střelba se automaticky zaměřuje na bowlingovou kouli ve scéně. To bylo zvoleno proto, aby hráč mohl využít obě ruce na klávesnici pro pohyb a další schopnosti. Výpočet směru projektilu využívá normalizovaný vektor mezi pozicí hráče a pozicí koule.

Vzorec pro výpočet směru je \texttt{(ballPosition - playerPosition).normalized()}. Tento vektor udává směr od hráče k cíli a jeho normalizace zajistí, že projektil má konstantní rychlost bez ohledu na vzdálenost.

Projektil se instancuje na pozici \texttt{firePoint}, což je bod na hráči odkud vylétá projektil. Rotace projektilu se nastaví podle směru letu pomocí \texttt{Mathf.Atan2} -- tento úhel se převede na stupně a nastaví se jako rotace kolem osy Z. Toto zajišťuje, že projektil je vizuálně orientován ve směru pohybu.

\begin{lstlisting}[language=csharp,caption={Listing 12: Metoda Shoot pro vytvoření projektilu}]
void Shoot()
{
    Vector2 direction = (enemyTransform.position - firePoint.position).normalized;

    GameObject bullet = Instantiate(bulletPrefab, firePoint.position, Quaternion.identity);
    Rigidbody2D rb = bullet.GetComponent<Rigidbody2D>();
    rb.linearVelocity = direction * bulletSpeed;

    float rotZ = Mathf.Atan2(direction.y, direction.x) * Mathf.Rad2Deg;
    bullet.transform.rotation = Quaternion.Euler(0, 0, rotZ);
}
\end{lstlisting}

\subsection{Vytvoření projektilu a nastavení fyzikálních parametrů}

Při výstřelu se volá \texttt{Instantiate()} s prefabem projektilu. Tato metoda vytvoří novou instanci projektilu ve scéně na specified position a rotation. Následně se získá komponenta \texttt{Rigidbody2D} na nově vytvořeném objektu a nastaví se rychlost pomocí \texttt{velocity = direction * bulletSpeed}.

Projektil se nikdy nemaže manuálně po čase. Vždy se zničí při kolizi s prostředím nebo koulí. Toto je implementováno ve skriptu \texttt{BulletStripeShot.cs}, který detekuje kolize a volá \texttt{Destroy(gameObject)} při kontaktu. Tento přístup šetří paměť a výkon, protože projektily, které minou cíl, nikdy nevisí ve scéně.
\newpage
\subsubsection{Omezený počet nábojů a systém přebíjení}

Hráč má k dispozici pouze 2 náboje, což je záměrný designový prvek. Střelba v naší hře není primárním prostředkem eliminace nepřátel, ale poslední možností záchrany. Proměnná \texttt{maxShots} určuje kapacitu a \texttt{reloadTime} dobu přebíjení.

Po vystřelení posledního náboje se spustí korutina \texttt{Reload()}, která čeká nastavenou dobu a následně obnoví počet nábojů na maximum. Během přebíjení nelze střílet, což vytváří strategický prvek.

Tento systém podporuje myšlenku, že výstřel má být pouze poslední možností záchrany a hráč se musí soustředit na zlepšování svých dovedností v pohybu a skocích. Omezená munice nutí hráče strategicky rozhodovat, kdy použít střelbu.
\subsection{Výstřel a kontrola podmínek}

Metoda \texttt{TryShoot()} se volá z \texttt{PlayerController} při inputu střelby. Kontroluje se několik podmínek: zda je střelba odemčená (volání \texttt{IsUnlocked()}), zda neprobíhá přebíjení \newline (proměnná \texttt{isReloading}) a zda jsou k dispozici náboje (\texttt{currentShots > 0}).

Pokud všechny podmínky platí, provede se výstřel voláním \texttt{Shoot()} a sníží se počet nábojů o jeden. Pokud byl vystřelen poslední náboj, spustí se přebíjení. Pokud některá z podmínek není splněna, metoda se ukončí bez akce.

Registrace inputu probíhá podobně jako u ostatních schopností. V \texttt{OnEnable()} se připojí callback na akci střelby a v \texttt{OnDisable()} se odpojí.
\subsubsection{Respawn a obnovení stavu střelby}

Při respawnu hráče se volá metoda \texttt{OnRespawn()}. Ta zastaví všechny korutiny (včetně přebíjení) a obnoví plný počet nábojů. Toto zajišťuje, že hráč má po každém respawnu čerstvou zásobu střeliva a přebíjení se nepřenáší mezi respawny.

Zastavení korutin je důležité, protože pokud by přebíjení běželo během respawnu, hráč by mohl přijít o munici, kterou by měl mít k dispozici. Vynulováním stavu se zajistí konzistentní herní podmínky po každém respawnu.

Metoda \texttt{OnRespawn()} je volána z \texttt{PlayerDeath} skriptu, který spravuje logiku smrti a respawnu hráče. Toto propojení zajišťuje, že všechny systémy hráče jsou správně resetovány při každém novém životě.
\newpage
\subsubsection{Chování projektilu po vystřelení}

Skript \texttt{BulletStripeShot.cs} se nachází na instanci projektilu a aktivuje se až po \newline vystřelení. V \texttt{OnTriggerEnter2D()} se kontroluje, s jakým objektem projektil kolidoval.

Při vstupu do trigger collideru bowlingové koule (která má tag \texttt{Ball}) se volá metoda \texttt{Stun()} na kouli. To způsobí, že koule se zpomalí a poskytne hráči čas na únik. Na rozdíl od tradičních her, kde střelba eliminuje nepřítele, v naší hře střelba pouze dočasně neutralizuje hrozbu.

Při srážce s prostředím (Collider2D na ground vrstvě) se projektil okamžitě zničí voláním \texttt{Destroy(gameObject)}. Toto je implementováno pomocí layer comparison nebo tag checking, aby se rozlišilo mezi prostředím a jinými objekty.

\begin{lstlisting}[language=csharp,caption={Listing 13: Skript BulletStripeShot}]
public class BulletStripeShot : MonoBehaviour
{
    public bool stun = true;
    public LayerMask groundLayer;

    private void OnTriggerEnter2D(Collider2D collision)
    {
        if (collision.CompareTag("Ball"))
        {
            BallMovement ball = collision.GetComponent<BallMovement>();
            if (ball != null)
                ball.Stun(stun);

            Destroy(gameObject);
        }
    }

    private void OnCollisionEnter2D(Collision2D collision)
    {
        if (((1 << collision.gameObject.layer) & groundLayer) != 0)
            Destroy(gameObject);
    }
}
\end{lstlisting}


\subsection{Bowlingová koule}

Hlavním nepřátelským objektem ve hře je bowlingová koule, která představuje konstantní hrozbu pro hráče. Koule se neustále pohybuje směrem dopředu a jejím úkolem je donutit hráče neustále se pohybovat vpřed. V této sekci bude podrobně popsána implementace chování bowlingové koule, včetně systému stunu a dynamické změny rychlosti.

\subsection{Přehled a základní struktura}

Hlavní nepřátelský objekt ve hře, bowlingová koule, je řízen skriptem \texttt{BallMovement.cs}. Ten je přidělen na herní objekt koule a stará se o její konstantní pohyb doprava, detekci stunu a dynamický růst rychlosti. Skript vyžaduje \texttt{Rigidbody2D} a \texttt{CircleCollider2D} přes atribut \texttt{[RequireComponent]}, což zajišťuje, že jsou vždy přítomny potřebné komponenty.

Koule představuje konstantní hrozbu, která nutí hráče neustále se pohybovat vpřed. Na rozdíl od tradičních nepřátel, koule nelze zničit, pouze dočasně zpomalit. Po uplynutí stunu se koule opět zrychlí na vyšší rychlost než předtím, což vytváří progresivní obtížnost v průběhu úrovně.
\newpage
\subsection{Pohyb a konstantní rychlost}

Koule se pohybuje konstantní rychlostí na ose X pomocí \texttt{Rigidbody2D}. V každém snímku se nastavuje \texttt{velocity.x = moveSpeed}, což zajišťuje konstantní pohyb doprava bez ohledu na kolize nebo jiné faktory. Tento přímý přístup k rychlosti je jednodušší a spolehlivější než použití sil.

Rigidbody je nastaven na \texttt{freezeRotation}, takže se koule neotáčí vlivem fyziky. Rotace animace koule se řídí samostatně v závislosti na rychlosti pohybu. Toto oddělení fyzikální rotace od vizuální animace umožňuje přesnější kontrolu nad prezentací.

\begin{lstlisting}[language=csharp,caption={Listing 14: Pohyb bowlingové koule}]
void Awake()
{
    rb = GetComponent<Rigidbody2D>();
    rb.freezeRotation = true;
}

void FixedUpdate()
{
    if (!_canMove) return;
    rb.linearVelocity = new Vector2(moveSpeed, rb.linearVelocity.y);
    if (animator != null)
        animator.speed = Mathf.Max(0.1f, rb.linearVelocity.magnitude / 5f);
}
\end{lstlisting}

\newpage
\subsection{Startovní prodleva a synchronizace s hráčem}

Před samotným pohybem běží korutina \texttt{StartMovingAfterDelay()}, která čeká \texttt{startDelay} sekund. Během této doby se \texttt{\_canMove} nastaví na false a pohyb se neprovádí. To dává hráči náskok před tím, než začne nebezpečí v podobě pronásledování koulí.

Tato prodleva je důležitá pro herní design, jelikož bez ní by se koule začala hned pohybovat ve směru hráče, což by mohlo působit nespravedlivě. Startovní delay poskytuje férový start a umožňuje hráči dostat se do bezpečnější vzdálenosti.

Stejná prodleva se restartuje při respawnu, aby po smrti měl hráč opět čas na přípravu. To zajišťuje konzistentní herní podmínky a zabraňuje situacím, kdy by se koule začala pohybovat k respawnujícímu se hráči okamžitě.
\newpage
\subsection{Stun systém a dočasné zpomalení}

Když projektil zasáhne kouli, volá se metoda \texttt{Stun()}. Ta inkrementuje counter \newline \texttt{\_activeStuns} a spustí korutinu \texttt{StunCooldown()}. Systém rozlišuje první stun a další stuny -- u prvního se uloží aktuální rychlost jako základ, u dalších se jen nastaví nízká rychlost.

Více stunů současně způsobuje, že rychlost zůstává na minimální hodnotě, dokud se všechny nestihnou vypořádat. Po uplynutí \texttt{stunDuration} se counter sníží a rychlost se obnoví na základní úroveň. Pokud je více stunů aktivních, rychlost zůstává na minimu.

Tento systém umožňuje hráči střílet více projektilů za sebou pro delší zpomalení, ale vzhledem k omezené munici (2 náboje) je toto omezeno. Navíc po každém stunu se základní rychlost zvýší, takže koule je postupně rychlejší.

\begin{lstlisting}[language=csharp,caption={Listing 15: Metoda Stun a StunCooldown}]
public void Stun(bool stun)
{
    if (!stun) return;
    _activeStuns++;
    StartCoroutine(StunCooldown());
}

private IEnumerator StunCooldown()
{
    if (_activeStuns == 1) _baseMoveSpeed = moveSpeed;
    moveSpeed = 5f;

    yield return new WaitForSeconds(stunDuration);

    _activeStuns = Mathf.Max(0, _activeStuns - 1);
    _baseMoveSpeed = Mathf.Min(_baseMoveSpeed + speedIncreasePerStun, maxSpeed);
    moveSpeed = _activeStuns == 0 ? _baseMoveSpeed : 5f;
}
\end{lstlisting}

\subsection{Růst rychlosti a progresivní obtížnost}

Důležité je, že \texttt{\_baseMoveSpeed} se ukládá a dále zvyšuje s každým dalším stunem. \newline Proměnná \texttt{\_baseMoveSpeed} se po každém stunu zvýší o \texttt{speedIncreasePerStun}, ale nikdy nepřesáhne \texttt{maxSpeed}. Pokud je více stunů aktivních současně, rychlost zůstává na minimu, dokud se všechny nestihnou vypořádat.

Tento systém vytváří progresivní obtížnost v průběhu úrovně. Každý úspěšný stun zpomalí kouli, ale po obnovení bude koule rychlejší než předtím. To nutí hráče vyhýbat se stunu, pokud není absolutně nutný, a soustředit se spíše na pohyb než na střelbu.

Maximální rychlost je omezena parametrem \texttt{maxSpeed}, aby koule nebyla příliš rychlá na to, aby ji hráč mohl uniknout. Správné nastavení tohoto parametru je klíčové pro vyvážení herní obtížnosti.
\subsection{Reset při smrti hráče}

Při smrti hráče se volá metoda \texttt{ResetStunAndSpeed()}. Ta zastaví všechny korutiny, vynuluje počet stunů a nastaví rychlost na předanou výchozí hodnotu. Zároveň restartuje startovací delay přes \texttt{StartMovingAfterDelay()}.

Tento reset zajišťuje, že po smrti hráče se koule vrátí do výchozího stavu. To je důležité pro férovost hry, protože hráč by neměl být penalizován za pokročilý stav koule z předchozího života. Po respawnu má hráč čas na únik a koule začíná opět od začátku.

Zastavení korutin je důležité, aby se předešlo chybám při restartu. Pokud by nějaká korutina běžela během resetu, mohla by způsobit nekonzistentní stav. Vynulování počtu stunů zajišťuje, že žádný stun nepřetrvává mezi životy.
\subsubsection{Animace rotace a dynamická rychlost}

Rychlost rotace animace koule se dynamicky mění podle aktuální rychlosti pohybu. Metoda \texttt{UpdateAnimatorSpeed()} počítá podíl aktuální rychlosti a základní rychlosti. \newline Výsledek se nastaví jako \texttt{animator.speed}, což ovlivňuje rychlost přehrávání animace.

Vzorec je \texttt{animator.speed = currentSpeed / baseSpeed}. Pokud je aktuální rychlost stejná jako základní, animator.speed je 1.0. Pokud je rychlost dvojnásobná, animator.speed je 2.0 a animace se přehrává dvakrát rychleji.

Toto zajišťuje, že vizuální rotace koule odpovídá skutečné rychlosti pohybu. Když koule stojí, rotace se zastaví. Když koule zpomalí (stun), rotace se zpomalí. Když koule zrychlí, rotace se zrychlí. Toto přispívá k celkovému vizuálnímu dojmu a konzistenci hry.

\subsection{Další skripty nepřítele}

Kromě základního pohybu bowlingové koule obsahuje nepřítel také další skripty, které rozšiřují jeho interakci s herním světem. Tyto skripty umožňují kouli ničit prostředí a reagovat na různé herní situace.
\newpage
\subsubsection{TileDestroy -- přehled a účel}

Skript \texttt{TileDestroy.cs} je přidělen na bowlingovou kouli a umožňuje jí ničit herní dlaždice v tilemap systému. Používá se k tomu Tilemap API, které umožňuje programově mazat dlaždice z úrovně. Tato mechanika přidává do hry dynamický prvek, kde prostředí může být poškozeno průchodem koule.

Tilemap systém v Unity ukládá dlaždice do gridu a umožňuje efektivní správu velkých herních světů. \texttt{TileDestroy} využívá \texttt{GetComponent<Tilemap>()} pro přístup k tilemap a \texttt{SetTile()} pro nastavení dlaždice na null (smazání).

Tato funkcionalita je klíčová pro herní design, protože vytváří časový tlak na hráče. Pokud koule dlouho stojí za hráčem, může zničit platformy a zablokovat cestu. Hráč musí neustále postupovat vpřed, aby si udržel cestu otevřenou.
\subsection{Tick-based ničení dlaždic a optimalizace}

Ničení dlaždic probíhá ve \texttt{FixedUpdate()} v taktovacích tickách, nikoli v každém snímku. Proměnná \texttt{timer} odpočítává čas mezi ticky, který je definován parametrem \texttt{destroyCooldown}. Jakmile timer dosáhne nuly, začne se procházet seznam kolizí a všechny dlaždice, jimiž koule prochází, se postupně odstraňují.

Důvod pro tick-based přístup je optimalizace. Kontrola dlaždic a jejich mazání je výpočetně náročná operace, zejména při vysoké rychlosti koule. Provedením této operace pouze jednou za několik snímků se výrazně sníží procesorová zátěž, aniž by to mělo viditelný dopad na herní zážitek.

Po dokončení ticku se timer resetuje na hodnotu \texttt{destroyCooldown} a cyklus pokračuje. Toto zajišťuje konzistentní chování bez ohledu na rychlost koule nebo snímkovou frekvenci.
\subsection{Limit dlaždic na tick a výkonová ochrana}

Aby se předešlo propadům výkonu při vysoké rychlosti koule, existuje parametr \texttt{maxTilesPerTick}. V každém ticku se zničí maximálně tento počet dlaždic. Jakmile je limit dosažen, další dlaždice se mažou až v dalším ticku. Parametr \texttt{destroyCooldown} určuje, jak často tick probíhá.

Toto omezení je důležité při vysoké rychlosti koule, kdy může projít stovkami dlaždic za sekundu. Bez limitu by to mohlo způsobit výrazné zpomalení hry. S limitem je zajištěno, že i při vysoké rychlosti zůstává výkon konzistentní.

Hodnota \texttt{maxTilesPerTick} se typicky nastavuje na 50-100 v závislosti na výkonu targetované platformy a typické rychlosti koule. Nižší hodnota = lepší výkon, vyšší hodnota = rychlejší ničení prostředí.
\newpage
\subsubsection{Bounds a WorldToCell -- převod souřadnic}

Pro každou kolizi se počítají \texttt{Bounds} a převádějí na souřadnice tilemap buněk pomocí \texttt{WorldToCell()}. \texttt{Bounds} representují ohraničující box collideru, zatímko \texttt{WorldToCell()} převádí world space souřadnice na grid space souřadnice tilemap.

Iteruje se přes všechny buňky v rozsahu bounds pomocí vnořených smyček. Pro každou buňku se zkontroluje, zda dlaždice existuje voláním \texttt{tilemap.HasTile(cell)}. Pokud dlaždice existuje, nastaví se na \texttt{null} voláním \texttt{tilemap.SetTile(cell, null)}.

Tento přístup zajišťuje, že všechny dlaždice v kontaktu s koulí jsou zničeny. Převod mezi world space a cell space je základní operací Tilemap API a umožňuje přesnou kontrolu nad ničením dlaždic.
\subsection{DestroyOnTouch -- přehled a funkčnost}

Skript \texttt{DestroyOnTouch.cs} řeší okamžité zničení objektů s tagem \texttt{Destroyable} při kontaktu s koulí. Na rozdíl od \texttt{TileDestroy}, který neničí dlaždice, ale celé herní objekty. Tento skript umožňuje kouli ničit různé objekty ve scéně jako jsou bedny, rozbitné předměty nebo dekorace.

Implementace využívá \texttt{OnTriggerEnter2D()} pro detekci vstupu do trigger collideru. Když koule vstoupí do triggeru jiného objektu, kontroluje se tag tohoto objektu. Pokud tag odpovídá \texttt{Destroyable}, objekt se zničí.

Přístup s triggery je vhodný pro objekty, které nepotřebují fyzikální kolizi. Trigger collidery detekují přítomnost bez fyzické interakce, což je užitečné pro objekty, které mají být zničeny, ale nemají ovlivňovat pohyb koule.
